% !PRINTED 128
\weberpage{128}
%<<CONT

sei, oder wenigstens keinen der anderen an Grösse übertrifft. Nach dem Satz 5., \S.\,38 lässt sich dann ein Werth $\alpha$ von $x$ so bestimmen, dass der absolute Werth $a$ von $f(\alpha)$ kleiner als $b_1$ wird, also:
\begin{equation}
 a<b_1\le b_2\le b_3\le\cdots\le b_{n-1}.
\tag{5}
\end{equation}

Wir begrenzen nun ein Gebiet $G$ für die Variable $x$ derart, dass \gesperrt{ausserhalb} dieses Gebietes der absolute Werth von $f(x)$ immer \gesperrt{grösser} als $a$ ist, so dass $G$ alle Punkte $x$ enthält, in denen $|f(x)|<a$ ist (aber auch noch andere Punkte). Dies ist nach \S.\,31, IV dadurch möglich, dass wir vom Nullpunkt als Mittelpunkt einen die Punkte $\alpha,\beta_1,\beta_2,
\ldots,\beta_{n-1}$ einschliessenden Kreis $(R)$ von hinlänglich grossem Radius $R$ legen und die Punkte $\beta_1,\beta_2,
\ldots,\beta_{n-1}$, in denen $|f(x)|$ ja grösser als $a$ ist, durch kreisförmige Hüllen von so kleinen aber nicht verschwindenden Radien $\varrho_1,\varrho_2,
\ldots,\varrho_{n-1}$ von diesem Kreise ausscheiden, dass im Inneren aller dieser Kreise $(\varrho_1),(\varrho_2),\ldots,(\varrho_{n-1})$ der absolute Werth $|f(x)|$ grösser als $a$ bleibt (\S.\,31, V). Dann umschliesst keiner dieser Kreise den Punkt $\alpha$, in dem $|f(x)|=a$ ist. Das von den Kreisen $(R),(\varrho_1),\ldots,(\varrho_{n-1})$ begrenzte zusammenhängende Flächenstück ist das Gebiet $G$.

% Exact 300-dpi crop from the registered Volume-I scan, printed Fig. 3.
\weberfigureimage{weber-b1-fig03.png}{0.58}


Wir bestimmen nun eine Zahlenreihe
\[
 \alpha,\alpha',\alpha'',\alpha''',\ldots
\]
derart, dass die absoluten Werthe von $f(\alpha),f(\alpha'),f(\alpha''),\ldots$, die wir mit
\[
 a,a',a'',a''',\ldots
\]
bezeichnen, immer abnehmen. Die so bestimmten Punkte $\alpha,\alpha',\alpha''$ bleiben alle im Inneren des Gebietes $G$, weil ja in ihnen $|f(x)|<a$ ist.

Wir bedienen uns dazu eines Verfahrens, ganz ähnlich dem zum Beweis des Satzes 5., \S.\,38 angewandten, nur dadurch vereinfacht, dass $f'(\alpha)$ schon von Null verschieden ist. Wir setzen, indem wir unter $\delta$ einen noch näher zu bestimmenden, positiven, echten Bruch verstehen, in der Entwickelung
%CONT>>
